% 启动流程

\section{内存布局}

这个 OS 的 \textit{Load Memory Address} (LMA) 如下图所示。这是它们被加载到硬件上时的物理内存位置，但并非虚拟内存位置。换句话说，虚拟内存位置 (VMA) 是给链接器看的，而 LMA 则是给加载器看的。

\newcount\foreachcount
\foreachcount=1
\begin{center}
\begin{tikzpicture}[node distance=0, every node/.style={minimum width=3.5cm, minimum height=1cm, align=center}]
  \foreach \name/\label/\ptr in {
    textlow/.text.low/0x8020'0000,
    ptroot/pt\_root/0x8020'1000,
    a0/a0 (hart id)/0x8020'2008,
    a1/a1 (FDT address)/0x8020'2010,
    text2/.text/0x8020'3000,
    rodata/.rodata/,
    data/.data/,
    bss/.bss/,
    stack/kernel stack/stack bottom
} {
    \ifnum\foreachcount=1
      \node (\name) [draw] {\label};
    \else
      \node (\name) [draw, below=of \prev] {\label};
    \fi

    % Text-only pointer node
    \IfEq{\ptr}{}{}{
      % if (ptr != "")
      \node (p\name) [minimum width=1cm, right=4cm of \name, anchor=south] {\ptr};
      \draw[->, thick] (p\name.west) -- ([xshift=1mm]\name.north east);
    }

    % Keep track of previous node
    \xdef\prev{\name}
    \global\advance\foreachcount by 1
  }

  % The stack top
  \node (last) [minimum width=1cm, right=4cm of stack, anchor=north] {. + 131072};
  \draw[->, thick] (last.west) -- ([xshift=1mm]stack.south east);
\end{tikzpicture}
\end{center}

其中，pt\_root 是页表的根。这个内存布局要求可以由链接器脚本 link.ld 完成。

从虚拟内存的角度而言，.text.low、pt\_root 以及 a0 与 a1 的 VMA 就等于它们的 LMA。其余部分的 VMA 都是 LMA 增加一个特定的偏移量 0xffff'ffc0'0000'0000 所得到的。这主要是因为 RISC-V 指令集\cite{rvPrivManual}的要求。

我采用的 Sv39 页表有 39 位的虚拟地址，但是放入的是 64 位寄存器。RISC-V 要求前 25 位必须 sign-extend，因此虚拟地址空间被强制分为了两部分。这个 OS 将会占据较高的半部分虚拟地址，而 user-space 程序将会使用较低的一半。上面的偏移量的前 26 位是 1，正是这个原因。

实际上，链接器——至少 GNU 的 ld——会将 VMA 按照 8-byte 对齐，但却不会对 LMA 这么做。这会导致一些让人摸不着头脑的 bug，因此我们需要在链接器脚本里通过 VMA 计算 LMA，而不是反过来。

\section{启动}

在启动时，我决定立即启用虚拟内存，而不先进行其他任何设置。甚至寄存器 sp 都没有被设置，因此栈也暂时无法使用。启动时的 a0 和 a1 是 hart id 和 FDT (\textit{flattened device tree}) 的地址，也需要顺手保存到上面图中的位置。

我先按照偏移量，开 16 个 1 GB 的页，将前 16 GB 物理内存映射至对应的虚拟地址。（QEMU 本身只有 128 MB 的物理内存，虽然理论而言这读了设备树才能知道。）接下来，我将第 8 GB 恒等映射，然后启用页表。这个恒等映射是必要的，不然启用页表后的下一条指令将会立即 page-fault。之后我会把它移除。

接下来，OS 将会跳转到真正的 .text 段开始执行。普通的 jal 是不可以的，因为指令集编码限制，它只支持 20 位的跳转，而我的跨度有 64 位那么大。同样地，la 也是不可以的——它实际上是 auipc，也只支持 20 位范围内的寻址。因此，我只能将高段的 .text 固定下来，让它留在 0x8020'3000 的位置。

接下来，我会启动一个 free-list allocator。它是一个物理分配器，而且只管理 16 MB 内存。这是为了让大部分 C++ 容器可以使用，请见 \cref{sec:container}。

下面我将启动 PLIC，接受硬件中断。不过这需要分为几步。

